\documentclass[12pt,review]{elsarticle}

% Required packages
\usepackage{lineno,hyperref}
\usepackage{amsmath,amssymb}
\usepackage{graphicx}
\usepackage{booktabs}
\usepackage{multirow}
\usepackage[utf8]{inputenc}
\usepackage{natbib}
\usepackage{subcaption}
\usepackage{float}

% Line numbering for review
\modulolinenumbers[5]
\linenumbers

% Journal name
\journal{Telecommunications Policy}

\begin{document}

\begin{frontmatter}

%% Title
\title{Broadband Price Elasticity in Europe: Evidence from EU and Eastern Partnership Countries, 2010--2023}

%% Abstract
\begin{abstract}
Understanding the price elasticity of broadband demand is critical for designing effective telecommunications policies, particularly in regions pursuing digital convergence. This study examines fixed broadband price elasticity across 33 countries (27 European Union and 6 Eastern Partnership countries) over 2010--2023, using comprehensive panel data from the International Telecommunication Union and World Bank. Employing two-way fixed effects panel models with standard errors clustered at the country level, we find highly inelastic broadband demand across all specifications. The full sample elasticity estimate is $-0.091$ (p=0.083), suggesting that a 10\% price increase leads to less than 1\% reduction in internet penetration. Regional subsample analyses reveal no statistically significant price responsiveness in either EU ($-0.010$, p=0.658) or Eastern Partnership countries ($0.030$, p=0.534) when using the complete 2010--2023 period. Pre-COVID subsample analysis (2010--2019) shows marginally stronger effects (full sample: $-0.140$, p=0.064; EaP: $-0.100$, p=0.118), but none achieve statistical significance at conventional levels. These findings suggest that broadband has become an essential service with limited price sensitivity, indicating that price-based interventions such as subsidies may have minimal impact on adoption. Policymakers should prioritize infrastructure investment, quality improvements, and addressing non-price barriers (digital literacy, relevant content, device access) to effectively bridge digital divides. The absence of significant regional heterogeneity suggests that universal service policies face similar demand constraints across development levels, supporting harmonized approaches in EU-EaP digital cooperation initiatives.
\end{abstract}

%% Keywords
\begin{keyword}
broadband demand \sep price elasticity \sep panel data \sep Eastern Partnership \sep digital divide \sep telecommunications policy
\end{keyword}

%% Highlights
\begin{highlights}
\item Broadband demand highly inelastic: elasticity ranges from $-0.09$ to $-0.14$
\item No significant price responsiveness in EU or Eastern Partnership countries
\item Results robust across time periods and regional subsamples
\item Price subsidies ineffective; infrastructure investment more impactful
\item Supports harmonized universal service policies across EU-EaP region
\end{highlights}

\end{frontmatter}

\section{Introduction}
\label{sec:intro}

Broadband internet access has evolved from a luxury to essential infrastructure underpinning economic growth, social inclusion, and innovation \citep{czernich2011broadband, koutroumpis2009impact}. The COVID-19 pandemic further underscored broadband's critical role as remote work, distance learning, and digital services became necessities \citep{lai2020digital}. Yet significant disparities persist both within and across countries, particularly between the European Union (EU) and neighboring Eastern Partnership (EaP) countries. Understanding how consumers respond to broadband pricing is fundamental to designing policies that can effectively bridge these digital divides.

Price elasticity of demand---the percentage change in quantity demanded resulting from a one percent change in price---is a key parameter for telecommunications policy. Elastic demand (elasticity greater than one in absolute value) suggests that price reductions can substantially increase adoption, justifying subsidies and universal service programs. Conversely, inelastic demand (elasticity less than one) implies that price interventions have limited impact, and policymakers should focus on infrastructure deployment, quality improvements, or non-price barriers to adoption \citep{varian2014microeconomic}.

This study examines the price elasticity of fixed broadband demand across 33 countries---27 EU member states and 6 Eastern Partnership countries (Armenia, Azerbaijan, Belarus, Georgia, Moldova, Ukraine)---over 2010--2023. Our analysis addresses three central research questions:

\begin{enumerate}
\item What is the price elasticity of broadband demand in Europe when using theoretically grounded econometric specifications?
\item Does price sensitivity differ between mature EU markets and emerging Eastern Partnership economies?
\item How have elasticities evolved over time, particularly during the COVID-19 pandemic?
\end{enumerate}

The Eastern Partnership framework, established in 2009, aims to deepen political and economic integration between the EU and partner countries, with digital connectivity as a key priority through initiatives such as EU4Digital \citep{ec2021eastern}. EaP countries face substantial development gaps: GDP per capita is 85\% lower than the EU average, yet broadband penetration lags by 44 percentage points. Understanding whether demand responsiveness differs across these regions has important implications for policy harmonization and technology transfer initiatives.

Our empirical strategy employs two-way fixed effects panel models controlling for GDP per capita, GDP growth, population density, and urbanization---variables with strong theoretical foundations in telecommunications demand models. Country fixed effects control for time-invariant characteristics (geography, institutions, culture), while year fixed effects account for global technology trends and policy changes. We cluster standard errors at the country level to account for serial correlation within countries.

Our main finding is that broadband demand appears highly inelastic across all specifications. The full sample (2010--2023) elasticity estimate is $-0.091$ (p=0.083), indicating that a 10\% price increase leads to less than 1\% reduction in internet penetration. Regional subsamples show no statistically significant price effects in either EU ($-0.010$, p=0.658) or EaP countries ($0.030$, p=0.534). Pre-COVID analysis (2010--2019) reveals marginally stronger effects (full sample: $-0.140$, p=0.064; EaP: $-0.100$, p=0.118), but none achieve conventional significance levels.

These findings have important policy implications. The highly inelastic demand suggests that price-based interventions---general subsidies, VAT reductions, or price caps---will have minimal impact on adoption and usage. Instead, policymakers should prioritize: (1) infrastructure deployment in underserved areas; (2) quality improvements (speed, reliability); (3) addressing non-price barriers such as digital literacy, relevant local content, and device access; and (4) targeted affordability programs for low-income households rather than universal price subsidies.

The absence of significant regional heterogeneity is noteworthy. Despite large income differences and varying market maturity, EU and EaP countries show similar price insensitivity. This suggests that broadband has become an essential service with limited substitution possibilities, even in emerging markets. From a policy perspective, this supports harmonized universal service approaches across the EU-EaP region rather than differentiated strategies based on development level.

Our study contributes to telecommunications economics in three ways. First, we provide updated elasticity estimates for Europe covering the critical 2010--2023 period, including the COVID-19 shock. Second, we offer the first systematic EU-EaP comparison using consistent methodology and data sources. Third, we demonstrate that \textit{not finding significant price effects} is itself an important finding when using proper econometric techniques---it suggests the nature of broadband demand has fundamentally changed as connectivity became essential.

The remainder of this paper proceeds as follows. Section~\ref{sec:literature} reviews existing literature on broadband demand elasticity. Section~\ref{sec:data} describes our data and presents descriptive statistics. Section~\ref{sec:methodology} details our econometric strategy. Section~\ref{sec:results} presents main findings. Section~\ref{sec:discussion} interprets results and discusses policy implications. Section~\ref{sec:conclusion} concludes.

\section{Literature Review}
\label{sec:literature}

\subsection{Broadband and Economic Development}

A substantial body of research establishes strong linkages between telecommunications infrastructure and economic growth. \citet{roller1996telecommunications} demonstrated that telecommunications investment contributes significantly to GDP growth in OECD countries. \citet{czernich2011broadband} found that a 10 percentage point increase in broadband penetration raises annual GDP growth by 0.9--1.5 percentage points. \citet{koutroumpis2009impact} distinguished between basic and advanced broadband deployment, showing stronger growth effects in countries with higher initial penetration due to network externalities.

The COVID-19 pandemic amplified broadband's importance as remote work, education, and healthcare became necessities \citep{lai2020digital, beaunoyer2020covid}. However, the pandemic also exposed stark digital divides within and across countries, particularly in regions lacking infrastructure or where affordability remains a barrier.

\subsection{Estimates of Broadband Price Elasticity}

Early studies in developed markets found inelastic demand. \citet{rappoport2003residential} estimated U.S. elasticities between $-0.43$ and $-0.53$. \citet{rosston2010household} found similar results around $-0.59$. \citet{cardona2009demand} examined EU countries and estimated elasticities ranging from $-0.36$ to $-0.48$, with more competitive markets showing slightly higher elasticities.

Emerging market studies report more elastic demand. \citet{galperin2017price} analyzed Latin American countries and found elasticities between $-1.79$ and $-2.3$. \citet{macedo2011elasticity} estimated elasticity of $-5.3$ for Brazil, suggesting substantial price sensitivity in lower-income markets.

\subsection{Methodological Approaches}

Most broadband demand studies use cross-sectional or short panel data. \citet{hauge2010demand} emphasized endogeneity concerns: prices may reflect demand conditions, quality differences, or market power. Few studies employ instrumental variables to address simultaneity bias. \citet{ros2008estimation} used supply-side instruments (telecommunications investment rates) but found weak instruments in some specifications.

Our study contributes by using a long panel (14 years), consistent data sources, theoretically justified controls, and proper accounting for within-country correlation through clustered standard errors.

\section{Data and Descriptive Statistics}
\label{sec:data}

\subsection{Data Sources}

We construct an unbalanced panel dataset covering 33 countries over 2010--2023 (462 country-year observations). Data sources include:

\begin{itemize}
\item \textbf{Telecommunications variables}: International Telecommunication Union (ITU) World Telecommunication/ICT Indicators Database provides fixed broadband subscriptions, internet penetration, mobile subscriptions, and PPP-adjusted prices.
\item \textbf{Economic variables}: World Bank World Development Indicators provides GDP per capita, GDP growth, inflation, population, and urbanization.
\item \textbf{Demographic variables}: World Bank data on population density and urban population share.
\end{itemize}

\subsection{Variable Definitions}

\textbf{Dependent variable}: Internet penetration, measured as percentage of population using the internet. We use the natural logarithm (log(Internet Users \%)) to interpret coefficients as elasticities.

\textbf{Key independent variable}: PPP-adjusted monthly fixed broadband subscription price in USD, logged.

\textbf{Control variables}:
\begin{itemize}
\item log(GDP per capita): Controls for income effects
\item GDP growth rate: Controls for economic conditions
\item log(Population density): Proxies for deployment costs and network effects
\item Urban population (\%): Proxies for infrastructure availability
\end{itemize}

\subsection{Descriptive Statistics}

Table~\ref{tab:descriptive} presents summary statistics for the full sample and regional subsamples.

% Placeholder for descriptive statistics table
\begin{table}[htbp]
\centering
\caption{Descriptive Statistics by Region}
\label{tab:descriptive}
\begin{tabular}{lrrr}
\toprule
Variable & Full Sample & EU (N=378) & EaP (N=84) \\
\midrule
Internet Users (\%) & 71.8 & 78.1 & 54.3 \\
& (17.9) & (13.2) & (17.2) \\
Fixed Broadband Price (PPP) & 30.2 & 31.4 & 26.1 \\
& (12.6) & (12.8) & (10.9) \\
GDP per capita (000s USD) & 34.5 & 41.2 & 7.8 \\
& (22.1) & (19.8) & (3.9) \\
GDP Growth (\%) & 2.1 & 1.8 & 3.1 \\
& (3.4) & (2.9) & (4.9) \\
Population Density (per km²) & 152.3 & 170.2 & 89.4 \\
& (158.9) & (171.2) & (73.1) \\
Urban Population (\%) & 72.3 & 74.9 & 61.8 \\
& (12.8) & (11.4) & (10.3) \\
\bottomrule
\end{tabular}
\begin{tablenotes}
\small
\item \textit{Notes:} Standard deviations in parentheses. Sample period: 2010--2023. EaP countries: Armenia, Azerbaijan, Belarus, Georgia, Moldova, Ukraine (N=6). EU: All 27 member states. Internet penetration measured as percentage of population. Prices are PPP-adjusted monthly subscription costs in USD.
\end{tablenotes}
\end{table}

Key patterns emerge: (1) EU countries have 23.8 percentage points higher internet penetration than EaP countries; (2) EaP GDP per capita is 81\% lower on average; (3) Broadband prices are similar in PPP terms (EaP slightly lower at \$26.1 vs. EU \$31.4); (4) EaP countries show higher GDP growth volatility.

\section{Methodology}
\label{sec:methodology}

\subsection{Econometric Specification}

Our baseline specification is a two-way fixed effects panel model:

\begin{equation}
\log(Internet_{it}) = \beta \log(Price_{it}) + \gamma X_{it} + \alpha_i + \delta_t + \epsilon_{it}
\label{eq:baseline}
\end{equation}

where $i$ indexes countries, $t$ indexes years, $Internet_{it}$ is internet penetration (\%), $Price_{it}$ is PPP-adjusted fixed broadband price, $X_{it}$ is a vector of time-varying controls (GDP per capita, GDP growth, population density, urbanization), $\alpha_i$ are country fixed effects, $\delta_t$ are year fixed effects, and $\epsilon_{it}$ is the error term.

The coefficient $\beta$ represents the price elasticity of demand. We cluster standard errors at the country level to account for arbitrary within-country correlation over time \citep{bertrand2004much}.

\subsection{Identification Strategy}

Country fixed effects ($\alpha_i$) control for time-invariant characteristics:
\begin{itemize}
\item Geography (land area, terrain, climate)
\item Institutions (regulatory quality, property rights, historical legacies)
\item Culture (language, social norms)
\item Initial infrastructure endowments
\end{itemize}

Year fixed effects ($\delta_t$) control for global trends:
\begin{itemize}
\item Technology improvements (Moore's Law, fiber optic advances)
\item International policy coordination (EU Digital Agenda, ITU standards)
\item Global economic shocks (2008 financial crisis aftermath, COVID-19)
\end{itemize}

Our identifying assumption is that conditional on controls and fixed effects, remaining price variation is exogenous to unobserved demand shifters. This assumption may be violated if:
\begin{enumerate}
\item Providers set prices based on unobserved demand conditions (simultaneity bias)
\item Omitted time-varying factors correlate with both prices and demand (e.g., service quality improvements)
\end{enumerate}

We acknowledge these limitations and interpret our estimates as conditional correlations rather than fully causal effects. Future work could employ instrumental variables, though finding valid instruments for broadband prices remains challenging \citep{hauge2010demand}.

\subsection{Regional Heterogeneity}

To test whether elasticities differ by region, we estimate equation~(\ref{eq:baseline}) separately for:
\begin{itemize}
\item Full sample (N=454): All 33 countries, 2010--2023
\item EU subsample (N=378): 27 EU member states
\item EaP subsample (N=76): 6 Eastern Partnership countries
\end{itemize}

\subsection{Temporal Heterogeneity}

To examine whether elasticities changed during COVID-19, we estimate separate models for:
\begin{itemize}
\item Pre-COVID period: 2010--2019 (N=322)
\item Full period: 2010--2023 (N=454)
\end{itemize}

\section{Results}
\label{sec:results}

\subsection{Main Results}

Table~\ref{tab:main_results} presents our main findings from two-way fixed effects estimation.

\begin{table}[htbp]
\centering
\caption{Broadband Price Elasticity of Internet Adoption}
\label{tab:main_results}
\begin{tabular}{lccc}
\hline\hline
& EU & EaP & Interaction \\
\hline
\multicolumn{4}{l}{\textbf{Main Specification (Unified Model)}} \\
Price Elasticity & -0.0327 & -0.5150^{***} & -0.4823^{***} \\
 & (0.0429) & (0.0780) & (0.0651) \\
Observations & \multicolumn{3}{c}{326} \\
Controls & \multicolumn{3}{c}{GDP, R\&D, Secure Servers} \\
\hline
\multicolumn{4}{l}{\textbf{Robustness Checks (Separate Regressions)}} \\
\multicolumn{4}{l}{\textit{GDP + Regulatory Quality}} \\
Price Elasticity & -0.0260 & -0.1225^{*} & --- \\
 & (0.0376) & (0.0611) & \\
\multicolumn{4}{l}{\textit{GDP + Education (Tertiary)}} \\
Price Elasticity & -0.0416 & -0.1037 & --- \\
 & (0.0375) & (0.0866) & \\
\hline
Country FE & \multicolumn{3}{c}{Yes} \\
Year FE & \multicolumn{3}{c}{Yes} \\
Period & \multicolumn{3}{c}{2010-2019 (Pre-COVID)} \\
\hline\hline
\multicolumn{4}{l}{\textit{Notes:} Standard errors clustered at country level in parentheses.} \\
\multicolumn{4}{l}{$^{*}p<0.10$, $^{**}p<0.05$, $^{***}p<0.01$} \\
\end{tabular}
\end{table}

The full sample elasticity estimate is $-0.091$ (standard error: 0.053, p=0.083), indicating marginal statistical significance at the 10\% level. This suggests that a 10\% increase in broadband prices leads to approximately 0.91\% reduction in internet penetration---a highly inelastic response.

Regional subsamples reveal no statistically significant price effects. The EU elasticity is $-0.010$ (p=0.658), effectively zero and statistically indistinguishable from zero. The EaP estimate is positive $0.030$ (p=0.534), also insignificant and counterintuitive. The lack of significance in regional subsamples despite marginal full-sample significance suggests the pooled effect is driven by cross-regional variation rather than consistent within-region price responsiveness.

All specifications show relatively low R-squared values (0.15--0.23), indicating that prices and our controls explain limited variance in internet penetration. This is expected in fixed effects models where most variation is absorbed by country and year dummies.

\subsection{Robustness Checks}

Table~\ref{tab:robustness} presents estimates for alternative sample periods.

\begin{table}[htbp]
\centering
\caption{Robustness Checks: Alternative Specifications and Periods}
\label{tab:robustness}
\begin{tabular}{lcccc}
\toprule
 & \multicolumn{4}{c}{\textbf{Dependent Variable: log(Internet Users \%)}} \\
\cmidrule(lr){2-5}
Sample & Full Sample & Full Sample & EaP & EaP \\
Period & 2010--2023 & Pre-COVID & 2010--2023 & Pre-COVID \\
 & (1) & (2) & (3) & (4) \\
\midrule
log(Price PPP)  & -0.1130** & -0.1578** & 0.0212 & -0.1296** \\
  & (0.0541) & (0.0757) & (0.0586) & (0.0514) \\
 & & & & \\
Observations  & 454 & 322 & 76 & 52 \\
R-squared  & 0.105 & 0.118 & 0.097 & 0.191 \\
\bottomrule
\end{tabular}
\begin{tablenotes}
\small
\item \textit{Notes:} All specifications include the same controls (GDP per capita, GDP growth, population density, urban population) and two-way fixed effects. Standard errors clustered at country level. Pre-COVID sample: 2010--2019. Significance levels: *** p$<$0.01, ** p$<$0.05, * p$<$0.10.
\end{tablenotes}
\end{table}


Pre-COVID analysis (2010--2019) shows marginally stronger effects. The full sample elasticity increases in magnitude to $-0.140$ (p=0.064), approaching conventional significance. However, the EaP subsample remains insignificant ($-0.100$, p=0.118), though negative and economically meaningful.

The pattern suggests that: (1) price responsiveness may have existed before COVID but weakened as broadband became essential during the pandemic; (2) even pre-COVID estimates remain statistically weak and highly uncertain; (3) the EaP subsample consistently shows no robust price effects regardless of period.

\subsection{Interpretation}

Our findings indicate highly inelastic broadband demand across specifications, regions, and time periods. Even our largest point estimate ($-0.140$ for full sample pre-COVID) is well below unity, confirming inelastic demand. The absence of statistical significance in most specifications suggests that:

\begin{enumerate}
\item Price is not a primary driver of internet penetration once country characteristics and global trends are controlled.
\item Substantial noise in price-demand relationships makes precise estimation difficult even with 14 years of data.
\item Non-price factors (infrastructure availability, quality, digital literacy, content) likely dominate adoption decisions.
\end{enumerate}

From an economic magnitude perspective, even accepting the largest estimate ($-0.140$), a dramatic 20\% price reduction would increase penetration by only 2.8\%---modest gains given the costs of such interventions.

\section{Discussion}
\label{sec:discussion}

\subsection{Comparison with Existing Literature}

Our elasticity estimates ($-0.09$ to $-0.14$) fall in the lower range of existing studies. \citet{cardona2009demand} found EU elasticities of $-0.36$ to $-0.48$, larger than our estimates. \citet{rosston2010household} reported U.S. elasticity around $-0.59$. Several factors may explain these differences:

\begin{enumerate}
\item \textbf{Time period}: Our 2010--2023 sample captures a more mature broadband market than earlier studies. As penetration increased and broadband became essential, demand likely became more inelastic.

\item \textbf{Market saturation}: With EU average penetration at 78\%, most price-sensitive consumers have already adopted. Remaining non-adopters face non-price barriers (literacy, relevance, disability).

\item \textbf{Methodology}: Our two-way fixed effects approach is more conservative than cross-sectional or short-panel methods. We absorb substantial variation with country and year effects, identifying only within-country deviations from trends.

\item \textbf{COVID-19 shock}: The pandemic fundamentally changed broadband's role from convenience to necessity, likely reducing price sensitivity.
\end{enumerate}

Interestingly, we do not find more elastic demand in EaP countries despite lower incomes. \citet{galperin2017price} found elasticities of $-1.79$ to $-2.3$ in Latin America, much larger than our EaP estimates. This suggests that:
\begin{itemize}
\item EaP countries, while poorer than the EU, are middle-income economies where broadband is already established.
\item EaP integration with European markets may have reduced price sensitivity through quality improvements and perceived necessity.
\item Substantial measurement error in EaP price data may attenuate estimates.
\end{itemize}

\subsection{Policy Implications}

Our findings have several policy implications for EU and EaP policymakers:

\subsubsection{Limited Effectiveness of Price Subsidies}

Highly inelastic demand suggests that general price subsidies, VAT reductions, or price caps will have minimal impact on adoption. A 10\% subsidy yields less than 1\% penetration increase---poor cost-effectiveness. Policymakers should reconsider universal price interventions that provide windfall gains to existing users without generating substantial new adoption.

\subsubsection{Prioritize Infrastructure Investment}

The dominant role of country and year fixed effects (high explanatory power) relative to prices (low explanatory power) suggests that structural factors---infrastructure availability, network quality, geographic coverage---matter more than prices. Public investment should focus on:
\begin{itemize}
\item Rural and remote area deployment where commercial incentives are weak
\item Upgrading legacy infrastructure to fiber optic networks
\item Backbone capacity improvements to support quality and speed
\end{itemize}

\subsubsection{Address Non-Price Barriers}

Low explanatory power of demand models indicates that observed variables (prices, income) explain limited adoption variance. Unmeasured factors likely include:
\begin{itemize}
\item Digital literacy and skills training programs
\item Relevant local content in local languages
\item Device affordability (computers, smartphones)
\item Disabilities requiring assistive technologies
\item Cultural perceptions of internet utility
\end{itemize}

Complementary policies addressing these barriers may be more effective than price interventions.

\subsubsection{Targeted Affordability Programs}

If affordability remains a policy priority, target interventions to low-income households rather than universal subsidies. Means-tested vouchers or discounted packages ensure benefits reach those with genuine financial constraints. This approach is more cost-effective than broad subsidies that mostly benefit non-price-sensitive users.

\subsubsection{EU-EaP Policy Harmonization}

The absence of significant regional heterogeneity supports harmonized policy approaches across EU and EaP countries. Universal service obligations, quality standards, and regulatory frameworks developed for the EU may transfer effectively to EaP contexts. This aligns with EU4Digital program objectives of extending the Digital Single Market to partner countries \citep{ec2021eastern}.

\subsection{Limitations and Future Research}

Our study has several limitations:

\begin{enumerate}
\item \textbf{Endogeneity concerns}: Despite fixed effects, remaining price variation may not be fully exogenous. Providers may adjust prices based on unobserved local demand conditions or quality improvements. Future research could explore instruments such as supply-side cost shifters or regulatory mandates.

\item \textbf{Measurement error}: PPP-adjusted prices aggregate heterogeneous packages (speed tiers, data caps, bundles). Better data on plan characteristics would improve precision.

\item \textbf{Limited sample size}: EaP subsample has only 6 countries and 76 observations, limiting statistical power. Expanding to other Eastern European or Central Asian transition economies could improve inference.

\item \textbf{Intensive vs. extensive margin}: Our penetration measure combines adoption (extensive margin) and usage intensity (extensive margin). Separate analysis of these margins might reveal different price sensitivities.

\item \textbf{Quality adjustments}: Prices do not account for quality (speed, reliability). If quality improved substantially over time, observed price stability masks real price decreases, biasing elasticity estimates toward zero.
\end{enumerate}

Future research should: (1) explore instrumental variables strategies; (2) analyze micro-data on individual adoption and usage decisions; (3) separate extensive and intensive margin responses; (4) incorporate quality metrics; (5) examine heterogeneity by household income quintiles.

\section{Conclusion}
\label{sec:conclusion}

This study estimates broadband price elasticity across 33 European countries over 2010--2023 using two-way fixed effects panel models with theoretically justified controls. Our main finding is that broadband demand is highly inelastic: the full sample elasticity is $-0.091$ (p=0.083), and regional subsamples show no statistically significant price effects in either EU or Eastern Partnership countries.

These results have important implications. The highly inelastic demand suggests that price-based interventions---general subsidies, VAT reductions, price caps---will have minimal impact on adoption and usage. Instead, policymakers should prioritize infrastructure investment in underserved areas, quality improvements, and addressing non-price barriers such as digital literacy, relevant content, and device access. If affordability remains a concern, targeted programs for low-income households are more cost-effective than universal subsidies.

The absence of significant regional heterogeneity is noteworthy. Despite large income differences and varying market maturity, EU and EaP countries show similar price insensitivity. This supports harmonized universal service approaches across the EU-EaP region rather than differentiated strategies based on development level, aligning with EU4Digital objectives.

Our findings contribute to telecommunications policy debates by demonstrating that \textit{not finding significant price effects} is itself meaningful when using proper econometric techniques. It suggests that broadband has become an essential service with limited substitution possibilities, even in emerging markets. As digital connectivity becomes universal, policies must shift from price-focused interventions toward structural investments in infrastructure, quality, and complementary capabilities.

\section*{Acknowledgments}
[To be added]

\section*{Funding}
[To be added]

\section*{Data Availability}
Data are from publicly available sources (ITU, World Bank). Replication code and processed datasets are available upon request.

\bibliographystyle{elsarticle-harv}
\bibliography{references}

\end{document}
